% =====================================================================
% Chapter 4 -- Dataset and Splits
% =====================================================================
\documentclass[../preamble.tex]{subfiles}
\begin{document}

\chapter{Dataset and Splits}
\label{ch:data}

\section{Source dataset: Figshare PCOS Ultrasound}
\label{sec:data-figshare}

The source-domain experiments use the public Figshare PCOS Ultrasound Dataset
\cite{figshare:pcos}.
The raw download contains more than 11{,}000 files arranged in infected and
non-infected directories. A byte-level MD5 audit found substantial duplication:
3{,}184 unique PCOS images among 6{,}784 infected files and 812 unique healthy
images among 5{,}000 non-infected files. The resulting unique dataset contains
3{,}996 images, with a PCOS-to-healthy ratio of $3.92{:}1$.

Deduplication precedes all split operations. If identical files are allowed to
cross a split boundary, a model can obtain an artificially high validation
score without learning a transferable feature. Exact-byte MD5 hashing detects
identical files; near duplicates arising from re-encoding or cropping are not
claimed to be removed.

\section{Target dataset: PCOSgen}
\label{sec:data-pcosgen}

The target-domain evaluation uses PCOSgen (Sundari et~al., 2025), with the
public release containing 4{,}668 images: 3{,}627 PCOS and 1{,}041 healthy.
The executed split contains a 3{,}200-image train pool and a separate
1{,}468-image test set. The train pool was divided into 2{,}560 training
images and 640 validation images. The held-out test set contains 549 PCOS
images and 919 healthy images.

The PCOSgen test set is not used for fine-tuning, HPO, epoch selection,
calibration fitting, or threshold selection. The results in Chapters~\ref{ch:results}
and~\ref{ch:discussion} therefore distinguish the fine-tuning pool from the
held-out external test.

\begin{table}[htbp]
  \centering
  \small
  \caption{Dataset composition across the executed pipeline.}
  \label{tab:dataset-composition}
  \begin{tabular}{lr@{\hspace{1em}}l}
    \toprule
    \textbf{Partition} & \textbf{Images} & \textbf{Notes} \\
    \midrule
    Figshare raw & 11,000+ & raw corpus before dedup \\
    Figshare PCOS & 3,184 & md5-unique infected \\
    Figshare non-PCOS & 812 & md5-unique non-infected \\
    Figshare unique & 3,996 & after md5 dedup; 3.92:1 imbalance \\
    PCOSgen raw & 4,668 & 3,627 PCOS / 1,041 healthy \\
    PCOSgen train pool & 3,200 & 2560 train + 640 val; 362 PCOS / 2838 healthy \\
    Fine-tune train & 2,560 & after stratified split \\
    Fine-tune val & 640 & best-epoch selection \\
    Held-out test & 1,468 & 549 PCOS / 919 healthy; never seen in training \\
    \bottomrule
  \end{tabular}
\end{table}


\section{Class imbalance}
\label{sec:data-imbalance}

Deduplication makes the source-domain class imbalance more visible because
the non-PCOS class contains more repeated files. The training loss uses
inverse-frequency weights:
\begin{equation}
  w_c = \frac{n_{\mathrm{total}}}{2n_c},
  \label{eq:class-weight}
\end{equation}
where $n_c$ is the number of examples in class $c$. A weighted sampler is
available as a separate data-level option; it is not silently combined with
loss weighting in the reported results unless explicitly selected in the
configuration. The target train pool is also imbalanced in the opposite
numerical direction: 362 PCOS versus 2{,}838 healthy images, approximately
$0.13{:}1$ PCOS-to-healthy. The held-out test set is less extreme at
$1{:}1.67$ PCOS-to-healthy.

\section{Split procedure}
\label{sec:data-split}

For the deduplicated Figshare data, the intended split is stratified 80/10/10:
10\% is first reserved for testing, then 10/90 of the remainder is reserved
for validation. The two-stage split preserves class proportions while leaving
approximately four hundred images in each evaluation partition.

Medical images can also leak patient identity if several frames from one
patient are distributed across partitions. The splitter therefore exposes a
group-aware fallback. When a patient key can be inferred and the number of
unique groups is smaller than the number of images, groups rather than images
are split. The Figshare implementation uses filename stems as a proxy when
metadata are absent. PCOSgen's released structure does not provide a verified
patient-level key, so this thesis does not claim patient-level separation for
that cohort; this is a limitation, not a result.

\section{Dataset routing}
\label{sec:data-routing}

The executed routing is:
\begin{enumerate}
  \item use Figshare to train the foundation models and assess source-domain
        performance;
  \item score foundation checkpoints zero-shot on the full PCOSgen corpus to
        expose domain shift;
  \item use only the PCOSgen train pool for adaptation and retain its 1{,}468
        test images as the final held-out evaluation;
  \item use the held-out predictions for final external metrics, confusion
        matrices, ROC/PR curves, and uncertainty summaries.
\end{enumerate}

This routing is why the zero-shot confusion matrix has $n=4{,}668$ while the
fine-tuned and ensemble confusion matrices have $n=1{,}468$. Those raw counts
must not be compared as if they came from the same population.

\section{Data limitations}
\label{sec:data-limitations}

The public datasets do not expose a complete demographic, acquisition-device,
or patient-identity metadata table. As a result, the thesis can measure
cross-dataset shift but cannot perform a meaningful subgroup fairness analysis
or guarantee patient-level independence in PCOSgen. Prospective multi-centre
data with verified patient IDs is required for clinical validation.

\end{document}
